% 实验二过程与结果
\subsection{实验过程与结果}

\subsubsection{实验设置}
构造 3 组 $\Lambda$-集合，每组包含 256 个明文，在矩阵坐标 $(1,0)$ 处遍历。分别测试 3 轮、4 轮以及无 MixColumns 情况下的异或和。
我们选取的明文如下：
% Plain 0
\[
P_0 = \begin{bmatrix}
123 & 34 & 115 & 178 \\
* & 195 & 84 & 74 \\
198 & 133 & 232 & 32 \\
225 & 167 & 181 & 207
\end{bmatrix}
\]

% Plain 1
\[
P_1 = \begin{bmatrix}
123 & 189 & 41 & 31 \\
* & 25 & 112 & 183 \\
202 & 38 & 211 & 215 \\
218 & 188 & 117 & 34
\end{bmatrix}
\]
% Plain 2
\[
P_2 = \begin{bmatrix}
48 & 16 & 194 & 55 \\
* & 237 & 85 & 225 \\
108 & 218 & 83 & 186 \\
154 & 71 & 135 & 139
\end{bmatrix}
\]
我们选择的密钥如下：
\[
key = 0xabf266f7a189b08e5db4a56ea585ae0f
\]
\subsubsection{结果记录}
根据实验运行结果，记录如下表：
(三组明文得到的结果都是一样的，对于4轮的情况，由于数据太长，只写了第一个)
\begin{table}[h]
    \centering
    \begin{tabular}{@{}cccc@{}}
        \toprule
        实验组别 & 轮数 & 包含 MixColumns & 异或结果 (XOR Sum) \\ \midrule
        1        & 3    & 是               & $[0, 0, 0]$ (平衡)         \\
        2        & 4    & 是               &$[213218948687531644049869954314427593478, ...]$ (不平衡) \\
        3        & 4-31 & \textbf{否}      & $[0, 0, 0]$ (始终平衡)     \\ \bottomrule
    \end{tabular}
    \caption{不同配置下的 AES 异或实验结果}
\end{table}
